\chapter{Terminology and Problem Description}
\label{chap:problem}

In this chapter, we will focus on the formal description of the extended 3D Bin Packing Problem.
To make this possible, we first need to introduce the terminology related to the representation of 3D space and the objects within it.



\section{Coordinates}

The \textit{coordinates} are defined as a triplet $(x, y, z)$, representing a position relative to the origin $(0, 0, 0)$, where:
\begin{itemize}
    \item $x$ denotes the position along the horizontal axis,
    \item $y$ denotes the position along the depth axis,
    \item $z$ denotes the position along the vertical axis.
\end{itemize}

Each coordinate value must be a non-negative integer:
\[
x, y, z \in \mathbb{Z}_{\geq 0}.
\]

The origin $\mathbf{0} = (0, 0, 0)$ serves as the reference point from which all coordinates are measured.



\section{Sizes}

The \textit{sizes} of a rectangular object are defined as a triplet $(X, Y, Z)$, where:

\begin{itemize}
    \item $X$ denotes the length of the edge along the horizontal axis,
    \item $Y$ denotes the length of the edge along the depth axis,
    \item $Z$ denotes the length of the edge along the vertical axis. 
\end{itemize}

Each of these values must be a non-negative integer:
\[
X, Y, Z \in \mathbb{Z}_{\geq 0}.
\]




\section{Region}

A \textit{region} is a rectangular object defined by a pair of coordinates $(\mathit{start}, \mathit{end})$, where:

\begin{itemize}
    \item $\mathit{start}$ denotes the coordinates of the vertex that is closest to the origin $(0, 0, 0)$,
    \item $\mathit{end}$ denotes the coordinates of the vertex that is furthest from the origin $(0, 0, 0)$.
\end{itemize}

It is required that the start coordinate is less than or equal to the end coordinate along each axis:
\[
\mathit{start}_x \leq \mathit{end}_x, \quad \text{and} \quad  
\mathit{start}_y \leq \mathit{end}_y, \quad \text{and} \quad 
\mathit{start}_z \leq \mathit{end}_z.
\]


\section{Rotation}



A \textit{rotation} of an item is represented by an element of the set of admissible rotations
\[
\mathcal{R} = \{ \rho_{XYZ}, \rho_{XZY}, \rho_{YXZ}, \rho_{YZX}, \rho_{ZXY}, \rho_{ZYX} \},
\]
where each rotation corresponds to a permutation of the dimensions:
\[
\begin{aligned}
\rho_{XYZ}(X, Y, Z) & = (X, Y, Z), \\
\rho_{XZY}(X, Y, Z) & = (X, Z, Y), \\
\rho_{YXZ}(X, Y, Z) & = (Y, X, Z), \\
\rho_{YZX}(X, Y, Z) & = (Y, Z, X), \\
\rho_{ZXY}(X, Y, Z) & = (Z, X, Y), \\
\rho_{ZYX}(X, Y, Z) & = (Z, Y, X).
\end{aligned}
\]

Equivalently, a rotation can be viewed as a function \(\rho \in \mathcal{R}\) that maps the size vector of an item to one of its six admissible permutations.

This approach maintains the values as non-negative integers but changes their assignment to axes, effectively rotating the item in discrete 90-degree steps.





\section{Relationships Between Regions}

Given two \textit{regions} $A = (A_\text{start}, A_\text{end})$  and $B = (B_\text{start}, B_\text{end})$, we define several spatial relationships.  

All comparisons between coordinates are interpreted componentwise, for instance, $A_\text{start} \leq B_\text{start}$ means:
\[
A_\text{start}^x \leq B_\text{start}^x \quad \text{and} \quad A_\text{start}^y \leq B_\text{start}^y \quad \text{and} \quad A_\text{start}^z \leq B_\text{start}^z.
\]

\subsection*{Subregion}
We say that region $A$ is a \emph{subregion} of region $B$ if:
\[
A_\text{start} \geq B_\text{start} \quad \text{and} \quad A_\text{end} \leq B_\text{end},
\]
meaning that $A$ lies entirely within the boundaries of $B$ in all three dimensions. We write $A \subseteq B$.

\subsection*{Overregion}
We say that region $A$ is an \emph{overregion} of region $B$ if $B$ is a subregion of $A$:
\[
A_\text{start} \leq B_\text{start} \quad \text{and} \quad A_\text{end} \geq B_\text{end}.
\]
We write $A \supseteq B$.


\subsection*{Intersection}
Two regions $A$ and $B$ \emph{intersect} if there is a non-empty volume that they share.

Formally, define open intervals for each dimension:
\[
I_A^d = (A_\text{start}^d, A_\text{end}^d), \quad I_B^d = (B_\text{start}^d, B_\text{end}^d), \quad \text{for } d \in \{x, y, z\}.
\]

The regions intersect if and only if the intersection of intervals is non-empty in all dimensions:
\[
A \cap B \neq \emptyset \iff \forall d \in \{x, y, z\}: I_A^d \cap I_B^d \neq \emptyset.
\]

This definition ensures that the overlapping region has strictly positive volume.



\subsection*{Volume Function}

We define a \emph{volume function} $\mathrm{V}$ that returns the volume of sizes or regions:

\begin{itemize}
    \item given sizes $D = (X, Y, Z)$, the volume is:
    \[
    \mathrm{V}(D) = X \cdot Y \cdot Z,
    \]
    
    \item given a region $R = (R_\text{start}, R_\text{end})$, the volume is defined as:
    \[
    \mathrm{V}(R) = (R_\text{end}^x - R_\text{start}^x) \cdot (R_\text{end}^y - R_\text{start}^y) \cdot (R_\text{end}^z - R_\text{start}^z).
    \]
\end{itemize}

The result of $\mathrm{V}$ is always a non-negative integer (possibly zero).













\section{Problem Description}

We are given a rectangular container with properties $C_\text{properties}$, consisting of sizes $S$ and a maximum allowed weight $W_\text{max}$, 
and a list $B$ of $N$ properties of rectangular input boxes. We call this pair the \emph{packing input}. 
Each input box $i$ with properties $B_i$, $i = 1, \dots, N$ has sizes $s_i$ and weight $w_i$.
We model the container as a region $(\mathbf{0}, S)$, 
where the sizes $S$ are interpreted as the end coordinate.
Similarly, when a box $i$ with properties $B_i$ is placed inside the container with a rotation $\rho$, 
it occupies a region $r_i = (p_i, p_i + \rho(s_i))$, where $p_i$ is the start coordinate of the input box within the container. 
We require that $r_i$ is a subregion of $(\mathbf{0}, S)$ as required by constraint (C2).

Let $\mathcal{B}_c$ denote the set of boxes placed in container $c$. 


We define the following:


\begin{itemize}
    \item the total weight of container $c$ is:
    \[
    W_c = \sum_{i \in \mathcal{B}_c} w_i,
    \]
    
    \item the total volume occupied in container $c$ is:
    \[
    \mathrm{V}(c) = \sum_{i \in \mathcal{B}_c} \mathrm{V}(r_i).
    \]
\end{itemize}

\subsection*{Constraints}

We impose the following constraints:

\begin{itemize}
    \item[\textbf{(C1)}] each box $i$ is placed in exactly one container,
    
    \item[\textbf{(C2)}] each box $i$ must lie entirely within its container $c$:
         \[r_i \subseteq (\mathbf{0}, S),\]
    
    \item[\textbf{(C3)}] no two boxes $i$ and $j$ within the same container $c$ may intersect:
    \[
    \forall i \neq j \in \mathcal{B}_c: \quad r_i \cap r_j = \emptyset,
    \]
    
    \item[\textbf{(C4)}] the total weight of any container $c$ does not exceed its weight limit:
    \[
    W_c \leq W_\text{max}
    \],
    
    \item[\textbf{(C5)}] boxes may only be rotated by 90° along any axis,
    
    \item[\textbf{(C6)}] every point of the base of a box must coincide either with the container floor or with the top surface of one or more other boxes.
\end{itemize}

\subsection*{Objective}

The goal is to assign all boxes to containers in a way that satisfies all constraints (C1)--(C6), while minimizing the number of used containers.
Compared to the classical 3D bin packing problem, this version introduces two key additional constraints: the maximum weight limit (C4), 
and the full-support condition (C6) that ensures physical stability of stacked boxes.




\section{Fitness Function}
\label{sec:fitness-function}

Let $\mathcal{C}$ be the set of used containers.
The objective of the fitness function $\mathrm{F}$ is to numerically represent the quality of the solution. 
The lower the fitness, the better the solution should be. 
The primary optimization goal is to minimize the number of containers used $|\mathcal{C}|$. 
In case two solutions use the same number of containers, the one that utilizes the space more efficiently—i.e., with a higher minimum fill ratio (volume or weight) is considered better.
Combining these two requirements, the fitness function is defined as:
\[
\mathrm{F} = |\mathcal{C}| + \min_{c \in \mathcal{C}} \left( \max \left( \frac{\mathrm{V}(c)}{\mathrm{V}(S)}, \frac{W_c}{W_\text{max}} \right) \right).
\]


\section{Data Representation}

\subsection{Input}

The packing input consists of a description of the input container properties and a list of properties of the input boxes. 
The container properties $C_\text{properties}$ are sizes $S = (X, Y, Z)$ and maximum weight $W_\text{max}$. 
Each box $i$ has properties $B_i$ consisting of a unique identifier, sizes $s_i$ and weight $w_i$.
The input is provided in JSON format, where all numerical values are expressed in the same units.

\subsection{Output}

The output is a JSON object containing an array of containers.

Each container includes:
\begin{itemize}
    \item its unique identifier,
    \item total weight,
    \item total occupied volume,
    \item a list of packed boxes.
\end{itemize}

Each box entry in the list of packed boxes includes:
\begin{itemize}
    \item the box identifier,
    \item applied rotation,
    \item the region it occupies.
\end{itemize}

All numerical values are expressed in the same unit system as the input.

\subsection{Datasets}
\label{sec:datasets}
For the experimental evaluation, a dataset consisting of 120 problem instances inspired by Martello, Silvano and Pisinger \cite{martello2000} was generated. 
The dataset is divided into four classes, each containing 30 instances. 
For every class, 10 instances were generated for each value of $N \in \{50, 100, 150\}$. The container sizes are identical for all instances and are defined as $X = Y = Z = 100$.

\begin{table}[htbp]
\centering
\begin{tabular}{|c|c|c|c|}
\hline
\textbf{Type} & \textbf{x} & \textbf{y} & \textbf{z} \\
\hline
1 & $[1, X/2]$ & $[Y/2, Y]$ & $[Z/2, Z]$ \\
2 & $[1, X/2]$ & $[1, Y/2]$ & $[Z/2, Z]$ \\
3 & $[1, X/2]$ & $[1, Y/2]$ & $[1, Z/2]$ \\
4 & $[1, X]$ & $[1, Y]$ & $[1, Z]$ \\
\hline
\end{tabular}
\caption{Box size limits for each type.}
\label{tab:box_types}
\end{table}

Classes 1–3 represent instances with a dominant geometric structure, while class 4 corresponds to completely general instances. 

For classes 1–3, the boxes are generated randomly with a probability of 60\,\% from the corresponding type and 20\,\% from each of the remaining two types. This approach preserves the main structural characteristics of each class while introducing additional variability.
The box size limits for each type are shown in Table~\ref{tab:box_types}. From these intervals, the sizes are selected uniformly.

In addition to the geometric properties, weights were assigned to the boxes due to the presence of a maximum weight constraint. Based on this, three different datasets were constructed. Each of them contains the same 120 problem instances, but the assigned weights differ.

In the original dataset, the boxes do not have weights. To simulate this, the weight of every box is set to 1, while 
the maximum allowed weight of each container is set to $X \cdot Y \cdot Z$, corresponding to the weight of a fully filled container with unit density. In practice, this means the weight constraint is never restrictive, as any container can accommodate all boxes by weight.

The remaining two datasets, referred to as the \emph{medium} and \emph{heavy} datasets, incorporate density. For each box, density is generated from a normal distribution with a specified mean density $\mu$ and standard deviation $0.1\mu$. The weight of each box is then computed as the product of the density and the box volume.

For the medium dataset, the mean density is set to $\mu = 1.2$, while for the heavy dataset it is set to $\mu = 1.5$.



